\documentclass{article}
\usepackage[utf8]{inputenc}
\usepackage[T1]{fontenc}
\usepackage[margin=2cm]{geometry}
\usepackage{hyperref}
\sloppy
\begin{document}
\section*{Historical source}
Encyclopaedia Britannica scan, source row 2446. Parse 3.5 extracted the text. Super 3.5 VL corrected it from an enlarged crop of each block. It can still contain recognition errors.
Page headers, footers, and page numbers are removed.

\subsection*{Untitled continuation (E0218)}
\textit{Continues an entry from an earlier page.}

be in powder when shaken, it is a sign the operation hath succeeded: but if it be in a cake, and doth not fall into powder on shaking the matras, it shews that your matter was not sufficiently roasted and dried in the iron pan during the preparatory operation."

Mr Lemeri hath shewn, that excrement is not the only matter capable of producing this phosphorus with alum; but that, on the contrary, almost all animal and even vegetable matters are fit for this combination; that though Mr Homberg mixed alum in equal quantities only with the fecal matter, it may be used in a much greater proportion; and, in certain cases, will succeed the better; that, according to the nature of the substances to be worked on, the quantity of that fat may be more or less increased; and that whatever is added, more than the dose requisite for each matter, serves only to lessen the virtue of the phosphorus, or even destroys it entirely; that the degree of fire applied must be different according to the nature of those matters; and, lastly, that salts containing exactly the same acid with that of alum, or the acid of those salts separated from its basis and reduced into spirit, do not answer in the present operation: which shews, says Mr Lemeri, that many sulphureous matters may be substituted for excrement in this operation; but that there are no salts, or very few, if any, that will succeed in the place of alum.

This phosophorus, made either by Mr Homberg's or by Mr Lemeri's method, shines both by day and by night. Besides emitting light, it takes fire soon after it is exposed to the air, and kindles all combustible matters with which it comes in contact; and this without being rubbed or heated.

Mess. Homberg and Lemeri have given the most probable and the most natural explanation of the cause of the accession and other phenomena of this phosphorus. What they say amounts in short to what follows.

Alum is known to be a neutral salt, confiting of the vitriolic acid and a calcareous earth. When this salt is calcined with the fecal matter, or other substances abounding in oil, the volatile principles of these substances, such as their phlegm, their salts, and their oils, exhale in the same manner as if they were distilled; and there is nothing left in the matras, when those principles are dissipated, but a charred matter, like that which is found in retorts wherein such mixtures have been decomposed by distillation.

This remainder therefore is nothing but a mixture of alum and charcoal. Now, as the acid of this salt, which is the vitriolic, hath a greater affinity with the phlogiston than with any other fulbiance, it will quit its basis to unite with the phlogiston of the coal, and be converted by that union into a sulphur. And this is the very case, of which we have certain proofs in the operation for preparing this phosphorus; for when, after the volatile principles of the oily matter are drawn off, the fire is increased, in order to combine closely together the fixed parts that remain in the matras, that is, the alum and the charred matter, we perceive at the mouth of the matras a small blue fulphoreous flame, and a pungent smell of burning sulphur. Nay, when the operation is finih-

ed, we find a real sulphur slicking in the neck of the ma-tras; and, while the phofhorus is burning, it hath plainly a frong sulphureous smell. It is therefore certain, that this phofhorus contains an adult sulphur; that is, a matter disposed to take fire with the greatest ease. But though sulphur be very inflammable, it never takes fire of itself, without being either in contact with some matter that is actually ignited, or else being exposed to a considerable degree of heat. Let us see then what may be the cause of its ascension, when it is a constituent part of this phofhorus.

We mentioned just now, that the acid of the alum quits its balls, in order to form a sulphur by combining with the phlogiston of the coal. This basis we know to be an earth capable of being converted into lime; and that it is actually converted into quick-lime by the calcination necessary to produce the phosphorus. We know that new made lime hath the property of uniting with water so readily, that it thereby contracts a very great degree of heat. Now when this phosphorus, which is partly constituted of the basis of the alum converted into quick-lime, is exposed to the air, the lime instantly attracts the moisture of which the air is always full, and by this means, probably, grows so hot as to fire the sulphur with which it is mixed. Perhaps also the acid of the alum is not totally changed into sulphur: some parts thereof may be only half-disengaged from its basis, and in that condition be capable of attracting strongly the humidity of the air, of growing very hot likewise by imbibing the moisture, and so of contributing to the accession of the phosphorus.

There is also room to think that all the phlogiston of the charred matter is not employed in the production of sulphur in this phosophorus, but that some part of it remains in the state of a true coal. The black colour of the unkindled phosophorus, and the red sparks it emits while burning, sufficiently prove this.

Pour some human urine into a glass alembic; set it in a water-bath, and distil till there remain only about a fortieth part of what you put in; or else evaporate the urine in a pan set in the balneum mariae till it be reduced to the same quantity. With this heat nothing will exhale but an insipid phlegm, smelling however like urine. The residuum will, as the evaporation advances, become of a darker and darker russet, and at last acquire an almost black colour. Mingle this residuum with thrice its weight of sand, and distil it in a retort set in a reverberating furnace, with the usual precautions. At first there will come over a little more insipid phlegm like the former. When the matter is almost dry, a volatile spirit will rise. After this spirit, white vapours will appear on increasing the fire; a yellow oily liquor will come off, trickling down in veins; and together with this liquor a concrete volatile salt, which will stick to the sides of the receiver. At last there will come over a deep-coloured fetid oil. In the retort there will remain a saline earthy residuum, which being lixiviated will yield some fetid salt.


\section*{Source record}
\begin{description}
\item[Scan ID] \texttt{britannica:309f765b1ecc83e0a7f7762e44b4a927a7555443:train:2446}
\item[Image file] \texttt{page-2446-692d680df6a9.jpg}
\item[Image SHA-256] \texttt{692d680df6a9023fab76d686360d380fd4e7491a6a607a6f10835cc740d21619}
\item[Dataset] NationalLibraryOfScotland/encyclopaedia\_britannica\_illustrated at revision \texttt{309f765b1ecc83e0a7f7762e44b4a927a7555443} (CC0-1.0)
\item[Extraction] model \texttt{nvidia/nemotron-3.5-parse}, raw reply \texttt{extract/raw/page-2446-692d680df6a9.json}
\item[Raw reply SHA-256] \texttt{0613d3822d9f2de9b3768b53098e4ee0042908908600fedf99265c8fa26bc99f}
\item[Correction] model \texttt{nvidia/Nemotron-3.5-Super-VL}, block crops, raw replies \texttt{correct/raw/crop\_nothink/page-2446-692d680df6a9-b*.json}; super\_vl: 10
\item[Running header] C H E M I S SI T R Y. 17
\end{description}
\end{document}
